\documentclass[11pt,a4paper]{article}
\usepackage[italian]{babel}
\usepackage[margin=2.5cm]{geometry}
\usepackage{parskip}
\usepackage{amsmath, amssymb, amsthm}
\usepackage{booktabs}
\usepackage{array}
\usepackage{xcolor}
\usepackage{needspace}
\usepackage{enumitem}
\usepackage{listings}
\usepackage{tikz}
\usepackage[most]{tcolorbox}
\usepackage[hidelinks]{hyperref}
\usetikzlibrary{decorations.pathreplacing, shapes.geometric, arrows.meta, positioning}

% Niente vedove/orfane: i paragrafi non lasciano righe isolate a fine/inizio pagina
\widowpenalty=10000
\clubpenalty=10000

% Elenchi più compatti
\setlist{itemsep=2pt, parsep=0pt, topsep=4pt, partopsep=0pt}

% --- Riquadri con bordo nero, senza numero (si spezzano tra due pagine) ---
% Uso: \begin{definizione} ... \end{definizione}
%      \begin{teorema}[Nome] ... \end{teorema}  (il nome è facoltativo)
\tcbset{nota/.style={enhanced jigsaw, breakable, colback=white,
  colframe=black, boxrule=0.5pt, sharp corners}}
\NewTColorBox{definizione}{}{nota,
  before upper={\textbf{Definizione.}\ }}
\NewTColorBox{teorema}{o}{nota,
  before upper={\textbf{Teorema\IfValueT{#1}{ (#1)}.}\ }}
\NewTColorBox{proposizione}{o}{nota,
  before upper={\textbf{Proposizione\IfValueT{#1}{ (#1)}.}\ }}
\NewTColorBox{proprieta}{o}{nota,
  before upper={\textbf{Proprietà\IfValueT{#1}{ (#1)}.}\ }}
\NewTColorBox{assioma}{o}{nota, boxrule=1pt,
  before upper={\textbf{Assioma\IfValueT{#1}{ (#1)}.}\ }}

% --- Ambienti semplici (senza riquadro) ---
% Il titolo sta in cima al blocco, anche se il contenuto inizia con un riquadro o
% con codice; e non resta mai da solo in fondo a una pagina.
% Uso: \begin{esempio}[Titolo facoltativo] ... \end{esempio}
\newcommand{\newrunin}[2]{%
  \NewDocumentEnvironment{#1}{o}{%
    \par\Needspace{4\baselineskip}\addvspace{\medskipamount}\noindent
    \textbf{#2}\IfValueT{##1}{ (##1)}\textbf{.}\ \ignorespaces}%
  {\par\addvspace{\medskipamount}}}
\newrunin{esempio}{Esempio}
\newrunin{esercizio}{Esercizio}
\newrunin{osservazione}{Osservazione}

% V e F nelle tavole di verità
\newcommand{\V}{\textsf{V}}
\newcommand{\F}{\textsf{F}}

% --- Codice C ---
% Uso: \begin{codice} ... \end{codice}      codice C numerato
%      \begin{comando} ... \end{comando}    comandi da terminale
%      \begin{risultato} ... \end{risultato} output di un programma
%      \lstinline|printf("ciao");|           codice dentro al testo
\lstdefinestyle{c}{
  language=C,
  basicstyle=\ttfamily\small,
  keywordstyle=\color{blue}\bfseries,
  commentstyle=\color{gray}\itshape,
  stringstyle=\color{red!70!black},
  numbers=left,
  numberstyle=\tiny\color{gray},
  numbersep=8pt,
  columns=flexible,
  keepspaces=true,
  breaklines=true,
  showstringspaces=false,
  tabsize=4,
  morekeywords={include, define},
  % lettere accentate nei commenti
  literate={à}{{\`a}}1 {è}{{\`e}}1 {é}{{\'e}}1 {ì}{{\`i}}1
           {ò}{{\`o}}1 {ù}{{\`u}}1 {È}{{\`E}}1 {À}{{\`A}}1
}
\lstset{style=c}
\newtcblisting{codice}{listing only, nota, left=6mm,
  listing options={style=c}}
\newtcblisting{comando}{listing only, nota,
  listing options={basicstyle=\ttfamily\small, numbers=none, language={},
    columns=flexible, keepspaces=true, breaklines=true}}
\newtcblisting{risultato}{listing only, enhanced jigsaw, breakable,
  colback=black!5, colframe=black, boxrule=0.5pt, sharp corners,
  listing options={basicstyle=\ttfamily\small, numbers=none, language={},
    columns=flexible, keepspaces=true, breaklines=true}}

% --- Diagrammi di flusso (simboli standard) ---
\tikzset{
  terminale/.style = {ellipse, draw, minimum width=2.5cm, minimum height=0.9cm},
  io/.style        = {trapezium, trapezium left angle=70, trapezium right angle=110,
                      draw, minimum width=2.5cm, minimum height=0.9cm},
  processo/.style  = {rectangle, draw, minimum width=2.5cm, minimum height=0.9cm},
  decisione/.style = {diamond, draw, aspect=2, minimum width=2.5cm},
  freccia/.style   = {thick, -{Stealth}}
}

\title{Il limite notevole $\sin x / x$}
\author{Marco Cerbella}
\date{Lezione 13 -- 5 novembre 2026}

\begin{document}
\maketitle

\section{Il limite notevole $\frac{\sin x}{x}$}

\begin{teorema}[Limite notevole]
\[
  \lim_{x\to 0}\frac{\sin x}{x} = 1 .
\]
\end{teorema}

Si tratta di una forma indeterminata $\frac{0}{0}$. La dimostrazione usa la circonferenza goniometrica e il teorema del confronto. Sia $0 < x < \frac{\pi}{2}$, $A = (1, 0)$, $P$ il punto della circonferenza che forma con l'asse $x$ l'angolo $x$ e $H$ la sua proiezione sull'asse delle ascisse. Si ha:
\begin{itemize}
  \item $\overline{PH} = \sin x$;
  \item l'arco $\widehat{PA}$ ha lunghezza $x$, perché è l'arco sotteso dall'angolo $x$ misurato in radianti (raggio $1$).
\end{itemize}

\paragraph{Passo 1: $\sin x < x$.} Il segmento $\overline{PH}$ è un cateto e $\overline{PA}$ è l'ipotenusa di un triangolo rettangolo, quindi $\overline{PH} < \overline{PA}$. Inoltre la distanza tra due punti (in questo caso $P$ e $A$) in linea retta è sempre più corta di una curva che li unisce: $\overline{PA} < \widehat{PA}$. Dunque
\[
  \sin x = \overline{PH} < \overline{PA} < \widehat{PA} = x .
\]

\paragraph{Passo 2: $\tan x > x$.} Sia $Q$ l'intersezione tra la retta per $O$ e $P$ e la retta tangente alla circonferenza in $A$: $\tan x = \overline{AQ}$ e $\sin x = \overline{PH}$. Il triangolo $OAQ$ contiene il settore circolare $OPA$, quindi
\[
  \text{area}(OAQ) > \text{area}(\widehat{OPA}) .
\]
L'area del triangolo è $\frac{1}{2}\,\overline{OA}\cdot\overline{AQ} = \frac{1}{2}\cdot 1\cdot\tan x$, e l'area del settore è $\frac{1}{2}\,\widehat{PA}\cdot\overline{OA}^{\,2} = \frac{1}{2}\cdot x\cdot 1$. Quindi
\[
  \frac{1}{2}\tan x > \frac{1}{2}\,x \quad\Longrightarrow\quad \tan x > x .
\]

\paragraph{Passo 3: confronto.} Dividendo $\sin x < x$ per $x > 0$:
\[
  \frac{\sin x}{x} < \frac{x}{x} = 1 .
\]
Da $\tan x > x$, cioè $\dfrac{\sin x}{\cos x} > x$, dividendo per $x$ si ottiene $\dfrac{\sin x}{x} > \cos x$. Quindi, in un intorno destro di $0$,
\[
  \cos x < \frac{\sin x}{x} < 1 .
\]
Per $x \to 0$ si ha $\cos x \to \cos 0 = 1$ e l'ultima funzione è la costante $1 \to 1$. Per il teorema del confronto,
\[
  \lim_{x\to 0^+}\frac{\sin x}{x} = 1 .
\]

\begin{osservazione}
Per $x \to 0^-$ basta osservare che $\dfrac{\sin(-x)}{-x} = \dfrac{\sin x}{x}$, cioè la funzione è pari; la slide svolge il caso $x > 0$. Il limite bilatero vale allora $1$.
\end{osservazione}

\end{document}
